% 相邻样本经同一算法得到相近输出，从而在任意测试点上损失接近。
\begin{tikzpicture}[
  data node/.style={book node,fill=FigureInput!10,minimum width=27mm,minimum height=12mm,inner sep=3pt},
  alg node/.style={book node,fill=FigureModel!10,minimum width=21mm,minimum height=10mm},
  out node/.style={book node,fill=FigureOutput!8,minimum width=25mm,minimum height=10mm},
  note/.style={font=\kaishu\scriptsize,text=BookMuted,align=center}
]
  \node[data node] (s) {$S=(z_1,\ldots,z_i,\ldots,z_n)$};
  \node[data node,below=14mm of s] (sp) {$S'=(z_1,\ldots,z_i',\ldots,z_n)$};
  \draw[FigureLoss,very thick,<->] (s) -- node[right,note,text=FigureLoss]{只替换\\一个观测} (sp);

  \node[alg node,right=15mm of s] (a1) {同一算法$A$};
  \node[alg node,right=15mm of sp] (a2) {同一算法$A$};
  \node[out node,right=15mm of a1] (w1) {$w=A(S)$};
  \node[out node,right=15mm of a2] (w2) {$w'=A(S')$};
  \draw[book arrow] (s)--(a1);
  \draw[book arrow] (sp)--(a2);
  \draw[book arrow] (a1)--(w1);
  \draw[book arrow] (a2)--(w2);
  \draw[BookBlue,very thick,<->] (w1)--node[right,note,text=BookBlue]{损失变化受控}(w2);

  \node[book node,fill=FigureLoss!9,below right=4mm and 12mm of w1] (test)
    {任意测试点$z$上的损失差$\leq\beta_n$};
  \draw[book arrow,dashed] (w1)--(test);
  \draw[book arrow,dashed] (w2)--(test);
\end{tikzpicture}
